\scanpage{26}
\mspage{iv}
the Italian school: to give an algebraic proof of the completeness of the characteristic linear system of suitable complete continuous sys-tems on a surface in characteristic 0 (cf. [13], [18], [33] and [36]).

It is my pleasure to acknowledge at this point the great encourage-ment and stimulation which I have received from Oscar Zariski, John Tate, and Alexander Grothendieck. In addition, I want to give credit to the many sources from which I have taken a great deal:

Ch. 1 and 2 owe a great deal to D. Hilbert [14],

\S1.2. was developed independently by C. Chevalley, N. Iwahori, and M. Nagata,

\S2.2. is largely a theory of J. Tits,

Ch. 3 was worked out by J. Tate and myself,

\S4.3. includes an example of M. Nagata

\S4.5. is a theorem of B. Kostant,

In Ch. 5 and 7, the whole approach to moduli via functors is due to A. Grothendieck,

\S5.3. and \S5.4. follow suggestions of A. Grothendieck,

Ch. 6 is almost entirely the work of A. Grothendieck.

\begin{flushright}Harvard University\\Oct. 27, 1963\end{flushright}
\mspage{25}
